\section{Excitation system --- \texttt{esac8b}}
\label{sec:esac8b}

\texttt{esac8b} is the PSLF implementation of the IEEE~421.5 type~AC8B excitation
system: a PID voltage regulator driving a brushless AC exciter with a rotating
rectifier. Figure~\ref{fig:esac8b} shows the block structure with the parameters used
here. The regulator gains are $K_{PR} = 100$, $K_{IR} = 90$ and $K_{DR} = 20$ with
$T_{DR} = 0.1$\,s; the amplifier is a unity gain with $T_A = 0$, so the regulator
output $V_R$ is limited directly to $0 \leq V_R \leq 20$\,pu. The exciter integrator
has $T_E = 0.4$\,s, and the feedback $V_{FE}$ collects self-excitation $K_E \cdot V_E$,
saturation $S_E(V_E) \cdot V_E$ and the armature-reaction term $K_D \cdot I_{fd}$
with $K_D = 1.19$. The rectifier regulation is $K_C = 0.17$, applied through the
standard $F_{EX}(I_N)$ characteristic with $I_N = K_C \cdot I_{fd} / V_E$. The
field-current limit $V_{FEMAX} = 11.46$\,pu is active; the terminal-voltage and speed
multipliers are off.

\begin{figure}[!ht]
\centering
\resizebox{\linewidth}{!}{\begin{tikzpicture}[x=1cm,y=1cm]
  % main path
  \node[blk,minimum width=16mm] (tr) at (1.0,0) {$\dfrac{1}{1 + s \cdot T_R}$};
  \node[sum] (s1) at (2.9,0) {};
  \node[blk] (pid) at (6.0,0) {$K_{PR} + \dfrac{K_{IR}}{s} + \dfrac{K_{DR} \cdot s}{1 + s \cdot T_{DR}}$};
  \node[lim,minimum width=16mm] (amp) at (9.5,0) {$\dfrac{K_A}{1 + s \cdot T_A}$};
  \node[sum] (s2) at (11.2,0) {};
  \node[lim,minimum width=14mm] (te) at (12.5,0) {$\dfrac{1}{s \cdot T_E}$};
  \node[blk,minimum width=10mm] (pi) at (15.4,0) {$\Pi$};
  % feedback blocks
  \node[blk,minimum width=24mm] (vfe) at (12.8,-2.2) {$K_E + S_E(V_E)$};
  \node[blk,minimum width=24mm] (kd)  at (12.8,-3.4) {$K_D$};
  \node[blk] (fx)  at (17.8,-1.7) {$F_{EX}(I_N)$,\quad $I_N = \dfrac{K_C \cdot I_{fd}}{V_E}$};
  % signals
  \draw[sig] (-0.6,0) node[left,lbl]{$U_t$} -- (tr);
  \draw[sig] (tr) -- node[above,lbl]{$V_C$} (s1);
  \node[lbl] at ($(s1)+(-0.22,-0.22)$) {$-$};
  \draw[sig] (2.9,0.9) node[above,lbl]{$V_{REF}$} -- (s1);
  \draw[sig] (2.9,-0.9) node[below,lbl]{$V_S$ (PSS)} -- (s1);
  \draw[sig] (s1) -- (pid);
  \draw[sig] (pid) -- (amp);
  \node[lbl,above=0.5mm of amp] {$V_{RMAX}$};
  \node[lbl,below=0.5mm of amp] {$V_{RMIN}$};
  \draw[sig] (amp) -- node[above,lbl]{$V_R$} (s2);
  \node[lbl] at ($(s2)+(-0.22,-0.22)$) {$-$};
  \draw[sig] (s2) -- (te);
  \node[lbl,below=0.5mm of te] {$V_{EMIN}$};
  \draw[sig] (te) -- node[above,lbl,pos=0.3]{$V_E$} (pi);
  \draw[sig] (pi) -- (16.6,0) node[right,lbl]{$E_{fd}$};
  % V_E tap
  \coordinate (tap) at (14.2,0);
  \fill (tap) circle (1.2pt);
  \draw[thick] (tap) -- (14.2,-2.6);
  \draw[sig] (14.2,-2.2) -- (vfe.east);
  \draw[sig] (14.2,-2.6) -- (17.4,-2.6) -- (17.4,-2.15);
  % F_EX into Pi
  \draw[sig] (fx.west) -- (15.4,-1.7) -- (pi.south);
  % I_fd
  \draw[sig] (20.2,-3.4) node[right,lbl]{$I_{fd}$} -- (kd.east);
  \fill (18.6,-3.4) circle (1.2pt);
  \draw[sig] (18.6,-3.4) -- (18.6,-2.15);
  % V_FE summation
  \draw[thick] (vfe.west) -- (11.2,-2.2);
  \draw[thick] (kd.west) -- (11.2,-3.4) -- (11.2,-2.2);
  \fill (11.2,-2.2) circle (1.2pt);
  \draw[sig] (11.2,-2.2) -- node[left,lbl,pos=0.2]{$V_{FE} \leq V_{FEMAX}$} (s2);
\end{tikzpicture}
}
\caption{\texttt{esac8b} (IEEE AC8B) block diagram. $V_S$ is the stabilizer output; $V_{FE}$ is the exciter field current, limited to $V_{FEMAX}$.}
\label{fig:esac8b}
\end{figure}

The settings are identical for both machine families, Table~\ref{tab:esac8b}. The
saturation points $S_E(4.8) = 0.0222$ and $S_E(3.0) = 0.0219$ describe an exciter
that is almost linear over its working range.

\begin{table}[!ht]
\centering
\begin{tabularx}{\linewidth}{@{}l X l r@{}}
\toprule
\textbf{Parameter} & \textbf{Description} & \multicolumn{1}{c}{\textbf{Unit}} & \multicolumn{1}{c}{\textbf{Value}} \\
\midrule
\texttt{tr} & Terminal-voltage transducer time constant $T_R$ & s & 0.01 \\
\texttt{kpr} & Regulator proportional gain $K_{PR}$ & pu & 100 \\
\texttt{kir} & Regulator integral gain $K_{IR}$ & pu/s & 90 \\
\texttt{kdr} & Regulator derivative gain $K_{DR}$ & pu$\cdot$s & 20 \\
\texttt{tdr} & Derivative-block time constant $T_{DR}$ & s & 0.1 \\
\texttt{vrmax} & Regulator output upper limit $V_{RMAX}$ & pu & 20 \\
\texttt{vrmin} & Regulator output lower limit $V_{RMIN}$ & pu & 0 \\
\texttt{ka} & Amplifier gain $K_A$ & pu & 1 \\
\texttt{ta} & Amplifier time constant $T_A$ & s & 0 \\
\texttt{te} & Exciter time constant $T_E$ & s & 0.4 \\
\texttt{vfemax} & Exciter field-current limit $V_{FEMAX}$ & pu & 11.46 \\
\texttt{vemin} & Minimum exciter voltage $V_{EMIN}$ & pu & 0 \\
\texttt{ke} & Exciter self-excitation constant $K_E$ & pu & 1 \\
\texttt{kc} & Rectifier loading factor $K_C$ & pu & 0.17 \\
\texttt{kd} & Exciter demagnetising factor $K_D$ & pu & 1.19 \\
\texttt{e1} & Exciter voltage at saturation point 1 $E_1$ & pu & 4.8 \\
\texttt{se1} & Saturation at $E_1$, $S_E(E_1)$ & -- & 0.0222 \\
\texttt{e2} & Exciter voltage at saturation point 2 $E_2$ & pu & 3 \\
\texttt{se2} & Saturation at $E_2$, $S_E(E_2)$ & -- & 0.0219 \\
\texttt{vtmult} & Terminal-voltage multiplier on regulator limits (0 = off) & -- & 0 \\
\texttt{spdmlt} & Speed multiplier on exciter output (0 = off) & -- & 0 \\
\bottomrule
\end{tabularx}

\caption{\texttt{esac8b} parameters (common to all nine units).}
\label{tab:esac8b}
\end{table}
